\documentclass[sigconf, 10pt, nonacm, anonymous]{acmart}  % for paper submission of EWSN2025
%\documentclass[sigconf, 9pt, nonacm]{acmart} % for the camera ready of EWSN2025

%%
%% \BibTeX command to typeset BibTeX logo in the docs
\AtBeginDocument{%
  \providecommand\BibTeX{{%
    Bib\TeX}}}

%%
%% Submission ID.
%% Use this when submitting an article to a sponsored event. You'll
%% receive a unique submission ID from the organizers
%% of the event, and this ID should be used as the parameter to this command.
%%\acmSubmissionID{123-A56-BU3}

%%
%% For managing citations, it is recommended to use bibliography
%% files in BibTeX format.
%%
%% You can then either use BibTeX with the ACM-Reference-Format style,
%% or BibLaTeX with the acmnumeric or acmauthoryear sytles, that include
%% support for advanced citation of software artefact from the
%% biblatex-software package, also separately available on CTAN.
%%
%% Look at the sample-*-biblatex.tex files for templates showcasing
%% the biblatex styles.
%%

%%
%% The majority of ACM publications use numbered citations and
%% references.  The command \citestyle{authoryear} switches to the
%% "author year" style.
%%
%% If you are preparing content for an event
%% sponsored by ACM SIGGRAPH, you must use the "author year" style of
%% citations and references.
%% Uncommenting
%% the next command will enable that style.
%%\citestyle{acmauthoryear}

\usepackage{balance}
\newcommand{\dieter}[1]{{\color{blue}[Dieter: #1]}}
\newcommand{\matthias}[1]{{\color{purple}[Matthias: #1]}}

\settopmatter{printacmref=false}
\settopmatter{printfolios=true}

%% Dieter 
\usepackage{enumitem}
\usepackage{tikz}
\tikzset{
  font={\fontsize{9pt}{12}\selectfont}}
\usetikzlibrary{plotmarks}
\usetikzlibrary{intersections,backgrounds, patterns}
\usetikzlibrary{patterns,shapes.arrows}
\usetikzlibrary{positioning,fit,backgrounds}
\usetikzlibrary{spy}
\usepackage{pgfplots}
\usepackage{pgfplotstable}
\pgfplotsset{compat=newest} %for double ylabels on the right and the left
\usepgflibrary{arrows}% for more options on arrows
\usepgfplotslibrary{groupplots,dateplot}
\pgfplotsset{every axis/.append style={
	scaled x ticks = false,
	label style={font=\footnotesize},
	tick label style={font=\footnotesize},
	tick scale binop=\times}
}
\pgfkeys{/pgf/number format/.cd,%additional code
1000 sep={},%additional code
}

%%
%% end of the preamble, start of the body of the document source.
\begin{document}

%%
%% The "title" command has an optional parameter,
%% allowing the author to define a "short title" to be used in page headers.
\title{Lost in Compression? Matching Signals from Compressed I/Q Data}


% %%
% %% By default, the full list of authors will be used in the page
% %% headers. Often, this list is too long, and will overlap
% %% other information printed in the page headers. This command allows
% %% the author to define a more concise list
% %% of authors' names for this purpose.
% \renewcommand{\shortauthors}{Trovato et al.}

\author{Matthias, Dieter, Yago, Gerome, Sofie}

%%
%% The abstract is a short summary of the work to be presented in the
%% article.
\begin{abstract}
Time Difference of Arrival (TDOA)-based localization and collaborative sensing applications require accurate signal matching across spatially distributed receivers. When the signal structure and modulation are unknown, traditional decoding-based methods fail, necessitating the transmission of raw I/Q data, which is often infeasible due to bandwidth constraints. This paper investigates compression techniques that retain matching capability while significantly reducing data volume. We evaluate classical lossy methods (sample resolution reduction and downsampling) alongside a deep learning approach using a Siamese network to learn compact embeddings. Using real-world Mode S radar signals, we assess matching accuracy and compression ratio. Results show that moderate downsampling improves performance by suppressing noise, and that learned embeddings enable compression ratios two orders of magnitude higher than classical approaches with minimal accuracy loss. These findings offer practical guidance for efficient TDOA localization and collaborative sensing in bandwidth-constrained environments.
\end{abstract}

%%
%% The code below is generated by the tool at http://dl.acm.org/ccs.cfm.
%% Please copy and paste the code instead of the example below.
%% This is required for the camera-ready version; and is optional for papers submitted for review. 
\begin{CCSXML}
<ccs2012>
   <concept>
       <concept_id>10010583.10010588.10003247.10003248</concept_id>
       <concept_desc>Hardware~Digital signal processing</concept_desc>
       <concept_significance>500</concept_significance>
       </concept>
   <concept>
       <concept_id>10010147.10010257.10010258.10010260.10003697</concept_id>
       <concept_desc>Computing methodologies~Cluster analysis</concept_desc>
       <concept_significance>500</concept_significance>
       </concept>
   <concept>
       <concept_id>10003752.10003809.10010031.10002975</concept_id>
       <concept_desc>Theory of computation~Data compression</concept_desc>
       <concept_significance>500</concept_significance>
       </concept>
   <concept>
       <concept_id>10010405.10010432.10010988</concept_id>
       <concept_desc>Applied computing~Telecommunications</concept_desc>
       <concept_significance>300</concept_significance>
       </concept>
 </ccs2012>
\end{CCSXML}

\ccsdesc[500]{Hardware~Digital signal processing}
\ccsdesc[500]{Computing methodologies~Cluster analysis}
\ccsdesc[500]{Theory of computation~Data compression}
\ccsdesc[300]{Applied computing~Telecommunications}

%%
%% Keywords. The author(s) should pick words that accurately describe
%% the work being presented. Separate the keywords with commas.
\keywords{TDOA,localization,IQ data,compression}
%% A "teaser" image appears between the author and affiliation
%% information and the body of the document, and typically spans the
%% page.
% \begin{teaserfigure}
%   \includegraphics[width=\textwidth]{sampleteaser}
%   \caption{Seattle Mariners at Spring Training, 2010.}
%   \Description{Enjoying the baseball game from the third-base
%   seats. Ichiro Suzuki preparing to bat.}
%   \label{fig:teaser}
% \end{teaserfigure}



%%
%% This command processes the author and affiliation and title
%% information and builds the first part of the formatted document.
\maketitle

\section{Introduction}

% Had to move it here, otherwise placed to far below
\begin{figure*}[t]
    \centering
    \includegraphics[width=0.75\linewidth]{figures/system model.pdf}
    \caption{System model of our reference localization application}
    \label{fig:system_model}
\end{figure*}

Time Difference of Arrival (TDOA)-based localization is critical in various domains. In aviation, multilateration is widely used for airport and en-route surveillance, including wide-area multilateration for upper airspace monitoring. TDOA is also employed to localize ground transmitters via satellites~\cite{hawkeye2024rfgeo} or high-altitude platforms such as drones~\cite{UAVLocalization} or balloons~\cite{MILCOM23}. Additionally, it supports tracking Internet-of-Things (IoT) devices~\cite{IOTloc}.

In conventional TDOA-based systems, receivers are capable of demodulating the incoming signal. This allows for highly efficient compression by reducing the received high-bandwidth signal to its decoded payload. Each receiver transmits this payload and a timestamp to a Central Data Processor (CDP), which then performs signal matching, i.e., grouping receptions from multiple receivers belonging to the same transmission, and calculates the transmitter's location based on timestamp differences.

However, in scenarios where the modulation and coding scheme of the signal are unknown, it is unclear what information should be transmitted to the CDP to enable accurate signal matching. Implementations with lossless compression, such as the compression used in Electrosense \cite{rajendran2017electrosense}, may not achieve the substantial data volume reduction required for practical deployment in resource-constrained TDOA systems analyzing unknown signals. At high signal bandwidths or large networks, the resulting data volume can still exceed typical data link capacities, making the transmission of raw or even losslessly compressed data infeasible. For instance, with a sample rate of 10~MHz and 32-bit samples, each sensor generates 320~Mbps of data, well beyond typical data link capacities, especially in larger networks or systems with limited data link capabilities.

The central question this paper addresses is: \textit{How can we produce a compressed representation of unknown signals that minimizes data volume while retaining sufficient information to support accurate signal matching?}

This challenge is similar to the well-studied problem of device fingerprinting \cite{devicefingerprinting}, where classifiers attempt to determine whether a signal originates from a specific source based on some abstract representation of the signal (the fingerprint). However, our goal differs: instead of identifying the source, we are interested in matching individual transmissions, independently of the source. This is effectively a form of packet fingerprinting. The task is to decide, based on abstract (compressed) signal representations, whether two sensors observed the same transmission event.

To answer this, we evaluate various methods for compressing in-phase quadrature (I/Q) signal data with respect to their suitability for TDOA localization. Using real-world signal data recorded at 1030~MHz, the uplink frequency used by secondary surveillance radar and collision avoidance systems in aviation, we focus on the localization of ground radar systems, e.g., from airborne or spaceborne receivers. These applications are particularly sensitive to data volume due to limited power and data link bandwidth.

The efficient signal matching framework presented here is not limited solely to the localization context. Applications involving collaborative sensing and processing by distributed nodes, such as collaborative distributed modulation classification \cite{2024_dyspan}, also necessitate a mechanism to ensure that all nodes are processing the same transmissions. The ability to accurately match signals across these distributed receivers, a key feature of our proposed solution, is therefore beneficial in a wider array of collaborative systems beyond the specific localization scenario explored.

The main contributions of this work are as follows:
\begin{itemize}
    \item We analyze how reducing sample resolution and downsampling affect the accuracy of signal matching.
    \item We propose and evaluate a deep learning-based Siamese network approach using real-world recordings.
    \item We compare the performance of different methods in terms of data reduction and signal matching accuracy.
\end{itemize}

\section{Background and Assumptions}

\subsection{System Model}

We consider a TDOA-based localization system commonly used in aviation and signal intelligence (SIGINT) applications. The system model, depicted in \autoref{fig:system_model}, consists of the following steps and components:

\begin{enumerate}
    \item A transmitter emits a radio signal from an unknown location.
    \item At least three receivers at different locations detect the signal and timestamp its arrival using synchronized clocks. Note that for 3D localization, at least four receivers are needed. However, we limit our analysis to the 2D case without loss of generality.
    \item Each receiver $R_i$ forwards signal information $S_i^j = (t_i, p_j)$ to a central data processor (CDP), where $t_i$ is the timestamp and $p_j$ is some information describing the signal, such as the decoded payload. Note that when referring to receptions by different receivers of a single transmission, we omit the subscript $j$.
    \item The CDP processes the received signal detection data in three steps:
    \begin{enumerate}[leftmargin=*]
        \item \textbf{Signal Matching:} The CDP groups receptions $S_i^j$ from different receivers that correspond to the same transmission, yielding sets of receptions $\{S_i^j\ |\ i = 1, \dots, N\}$ for each transmission $S^j$.  \label{step:signal_matching}
        \item \textbf{Location Estimation:} If $N \ge 3$, use either a closed-form solution ($N = 3$) or an iterative algorithm ($N > 3$) to estimate the transmitter’s 2D location based on timestamp differences. \label{step:location_estimation}
        \item \textbf{Tracking:} The CDP may additionally apply tracking methods such as track association and filtering to the estimated locations to filter erroneous localization estimates and improve localization accuracy across multiple measurements. \label{step:tracking}
    \end{enumerate}
\end{enumerate}
A summary of the system model is depicted in \autoref{fig:system_model}.
\bigskip

The computational complexity of signal matching (step~\ref{step:signal_matching}) significantly impacts the downstream localization and tracking steps. If signal rates are low, simple plausibility checks, such as TDOA constraints based on radio horizon, suffice for accurate matching. However, in high-rate environments, multiple candidate signals may exist within a valid TDOA window, increasing the risk of false matches. False matches not only degrade localization accuracy but also increase computational load in steps~\ref{step:location_estimation} and \ref{step:tracking}, as the CDP must process many invalid position estimates, which the tracker must subsequently filter out.

In this work, we focus on challenging scenarios with high signal rates and limited signal knowledge. These include both civil and military contexts, such as localizing drones using unknown control protocols in crowded ISM bands, or identifying unknown and unlawful transmitters by spectrum regulators. Furthermore, we assume constrained data link capacities between receivers and the CDP. Hence, our objective is to transmit \emph{minimal} signal information while still enabling accurate matching at the CDP.

Finally, we note two special conditions in the considered system model that allow for trade-offs in the applied methods. First, decoding of signals at the CDP is not required for localization purposes. This allows for use of lossy compression of the raw signal data. Second, the system can tolerate a certain level of false matches, but missed matches (false negatives) must be minimized. This will be reflected in the choice of the target performance metric in the subsequent experimental analysis.

\subsection{In-Phase and Quadrature Data}

In-phase and quadrature (I/Q) data is a common representation of received baseband signals, typically produced by software-defined radios. These I/Q samples are the output of analog-to-digital converters (ADCs) that digitize the down-converted signal, either in baseband or at an intermediate frequency. Each sample consists of an in-phase (I) and a quadrature (Q) component, usually represented as signed integers of fixed bit width. This bit width, referred to as the ADC resolution, affects the receiver's sensitivity and dynamic range. Typical commercial off-the-shelf (COTS) ADCs provide resolutions ranging from 8 to 16 bits.

The ADC's sample rate determines both the bandwidth limit of the receiver (according to the Nyquist-Shannon sampling theorem) and the resulting data volume. In our setup, we use a 12-bit ADC per I and Q channel, sampling at 12~Msamples per second. This yields a raw data rate of 288~Mbit/s per receiver. However, due to byte alignment in practical computer and network systems, 16 bits must be used to store each 12-bit value, increasing the actual data rate to 384~Mbit/s.

In theory, TDOA-based localization can be performed by transferring all raw I/Q data to the central data processor (CDP), which can then perform cross-correlation between signal streams from different receivers. A correlation peak indicates a match between signal receptions, and its lag gives the TDOA value. This method supports signal-agnostic localization but is rarely practical due to the high data bandwidth required, especially in scalable deployments.

For the remainder of the paper, we make the following additional assumptions:
\begin{itemize}
    \item Transmissions are of short duration, as typically observed in digital communication systems that use small packets. Long-duration transmissions such as analog voice or continuous wave jamming are out of scope.
    \item Each transmission is detected and timestamped at the receiver, for example using a squelch filter or more advanced signal detection methods~\cite{TODO}.
\end{itemize}
As a result, we only collect I/Q samples corresponding to actual transmissions. This alone provides a substantial reduction in data volume, which depends on the channel occupancy. For instance, if the channel is active 25\% of the time, the resulting data volume is reduced to approximately 25\% of the continuous I/Q stream.

\subsection{Correlation-based Signal Similarity}
\label{sec:similarity}

As part of the signal matching step at the CDP using traditional signal matching, we assess whether two receptions $S_i$ and $S_k$ correspond to the same transmission based on the similarity of their I/Q samples. Let $S_i[n]$ and $S_k[n]$ denote the $n$-th complex I/Q sample extracted from two receptions of a signal $S$. We then perform the following calculations to calculate the similarity of the two signals on a normalized scale from 0 (orthogonal) to 1 (identical).

To ensure comparability, both sequences are first normalized to a fixed length $N$ through truncation or zero-padding. To account for phase offsets, we then compute the phase difference $\Delta \phi$ as the median of the sample-wise differences between the phases of $S_i$ and $S_k$:
$$\Delta \phi = \text{median}_{1 \le n \le N}(\angle S_i[n] - \angle S_k[n])$$
The sequence $S_k$ is then phase-aligned with $S_i$:
$$S_k' = S_k \cdot e^{j \Delta \phi}$$
We define the similarity metric as the \emph{normalized} complex inner product:
\begin{equation}
\text{corr}(S_i, S_k) = \frac{\left|\sum_{n=1}^{N} S_i[n] \cdot S_k'[n]\right|}{\|S_i\| \cdot \|S_k'\|}
\label{eqn:iq_similarity}
\end{equation}
where $\|{\cdot}\|$ is the Euclidean norm of the complex vector.

We note that this metric is a complex-valued version of the \emph{cosine similarity}, adapted to I/Q signals. The phase alignment ensures that the global phase difference between two signals does not affect the similarity score. The result lies in the interval $[0, 1]$, where a value close to 1 indicates high similarity (identical up to phase and scaling), and values near 0 indicate dissimilar or approximately orthogonal signals.

If $\text{corr}(S_i, S_k) > T$ for a predefined threshold $T$, the CDP considers the two receptions as matching, i.e., receptions of the same transmission by different receivers. This correlation-based similarity check is applied during the Signal Matching step~\ref{step:signal_matching} of the CDP processing pipeline. Note that this thresholding rule introduces a trade-off between false positives and false negatives, the impact of which is explored in the evaluation section below.
\subsection{Deep Learning-based Compression and Matching using a Siamese Network}
\begin{figure}[t]
    \centering
    \includegraphics[width=\linewidth]{figures/png/siamese_model.png}
    \caption{Siamese network architecture}
    \label{fig:siamese_architecture}
\end{figure}

The deep learning approach leverages a compression model and the Siamese network architecture \cite{siamese} to generate an efficient signal representation for matching. In this architecture, as shown in figure \ref{fig:siamese_architecture}, each receiver employs an identical deep learning model to process detected radio signals. This model generates a compressed, lower-dimensional embedding that encapsulates the most essential features of the transmission. The use of a Siamese network, with its shared weights across identical branches (one conceptually residing at each receiver), ensures that the resulting embeddings from different receptions of the same signal reside in a common, comparable feature space. This allows for direct similarity assessment at the CDP. The resulting embedding vector, $\mathbf{z}$ (containing real-valued elements), represents a reduced data volume compared to the raw I/Q samples and is transmitted to the CDP.

At the CDP, the task of signal matching is performed by evaluating the similarity between the received embedding vectors. Given two embeddings, $\mathbf{z}_i$ and $\mathbf{z}_k$, corresponding to receptions $S_i$ and $S_k$, their similarity is inversely proportional to the Euclidean distance:
\begin{equation}
d(\mathbf{z}_i, \mathbf{z}_k)= \|\mathbf{z}_i-\mathbf{z}_k\| = \sqrt{\sum_{m=1}^{M} (z_{i,m} - z_{k,m})^2}
\label{eqn:euclidean_distance}
\end{equation}
where $M$ is the dimensionality of the embedding vector. A smaller Euclidean distance between two embeddings suggests a higher degree of similarity between the original received signals, indicating a greater likelihood that they originate from the same transmission.

Similar to the thresholding applied to the correlation score in the classical method, a predefined threshold $T$ will be applied to the Euclidean distance. If $d(\mathbf{z}_i, \mathbf{z}_k) < T$, the CDP considers the two receptions as a match. This threshold $T$ on the embedding distance will also introduce a trade-off between false positives and false negatives in the matching process, which will be analyzed in the evaluation section.
\subsection{Evaluation Methodology and Experimental Setup}

To evaluate our approach using real-world signal data, we focus on the practical application of locating secondary surveillance radar Mode S systems. Mode S radars emit interrogation signals directed at aircraft transponders using Differential Binary Phase-Shift Keying (DBPSK) modulation. Each interrogation signal has a duration of either 19.75~µs or 33.75~µs, corresponding to different Mode S uplink formats.

For this work, signal data was collected using a GRX3X receiver developed by SeRo Systems\footnote{\url{https://sero-systems.de}}. This receiver is capable of detecting and demodulating Mode S interrogations and provides a GPS-synchronized timestamp with nanosecond resolution for each detection. In addition to detected signal metadata, the GRX3X also outputs a continuous stream of time-synchronized I/Q samples at a rate of 12~MHz and an ADC resolution of 12 bits.

Recordings were conducted at a site approximately 2~km line-of-sight from Frankfurt International Airport. Due to the high traffic density and the presence of ground-based Mode S interrogators within direct view of the receiver, the environment represents a particularly challenging scenario with very high signal rates and transmissions that are often repeated within short time intervals.

\subsubsection{Signal Preparation and Ground Truth}

For each detected Mode S signal, we extract the corresponding block of I/Q samples from the raw data stream, beginning with the preamble pulses and ending after the DBPSK modulation block. To enable comparison across both short and long interrogations, all signals are normalized to a fixed window size of 512 samples through zero-padding or truncation, as described in \autoref{sec:similarity}. We note that this window size accommodates both short interrogations (237 samples) and long ones (405 samples).

In our evaluations throughout the paper, two signals are considered a true match in the CDP if the Hamming distance between their decoded symbol sequences is zero, i.e., they carry the exact same bit sequence. This strict criterion defines our ground truth for evaluating matching performance: only signals with identical payloads should be matched.

\subsubsection{Data Collection and Evaluation Methodology}

For our evaluation, we collected three datasets, each comprising 5,000 individual signal detections. Two datasets were used for training and validation of the Siamese neural network, while the third was reserved for evaluation. To ensure independence between sets, we ensured that the sets are disjoint, i.e., no signal in one dataset has a Hamming distance of zero to any signal in another set.

To assess the accuracy of our signal matching method, we performed pairwise comparisons between all 5,000 signals in the evaluation dataset, resulting in 12,497,500 unique signal pairs. For each pair, we applied the proposed compression techniques to the signal data and then tested whether they would still be matched correctly. The decoded bit sequences from the GRX3X receiver served as ground truth for determining true and false matches.

Using this ground truth, we compute standard evaluation metrics including precision, recall, and the F1 score, which will be discussed in the following section.

\subsubsection{Performance Metrics}
\label{sec:performance_metrics}

The goal of our approach is to minimize the data volume collected at the CDP while preserving high matching accuracy. Consequently, the two primary performance metrics considered are the \emph{compression ratio} and the \emph{matching accuracy}.

Matching accuracy is evaluated in terms of \emph{precision}, \emph{recall}, and their harmonic mean, the \emph{F1 score}. A precision of $1$ indicates no false matches, whereas a recall of $1$ implies that all true matches were correctly detected. Given the trade-off between false matches and missed matches introduced by lossy compression, the F1 score serves as a balanced measure of matching performance.

We note that the overall \emph{accuracy}, defined as the ratio of correctly classified pairs to all evaluated pairs, is not informative in this context, as non-matches dominate the dataset (99.7\%). Due to the very high true negative rate, the accuracy metric is largely determined by the abundance of true negatives and therefore does not meaningfully reflect the system's ability to detect true matches.

To measure the compression ratio, we use the full, uncompressed I/Q data of the detected transmissions as the baseline. This reflects the maximum amount of signal information available in our system. Given the 12~MHz sample rate, 12-bit ADC resolution (packed into 16-bit integers), and two channels (I and Q), the uncompressed data volume $V(S)$ for a signal $S$ of duration $\tau(S)$ is
$$V(S) = \tau(S) \times 12~\text{MHz} \times 2 \times 16~\text{bit} = \tau(S) \times 384~\text{Mbit/s}$$

Let $\tilde{V}(S)$ denote the size of the compressed representation of the detected signal $S$. The compression ratio $R$ for a given method is then determined by calculating the ratio between the total volume of the uncompressed data to the total volume of the compressed data over all $N$ detected signals $S_i$ in the evaluation dataset:
\begin{equation}
R = \frac{\sum_{i=1}^{N} V(S_i)}{\sum_{i=1}^{N} \tilde{V}(S_i)}
\label{eqn:compression_rate}
\end{equation}


Based on these metrics, we evaluate the impact of different compression methods and compression levels on matching accuracy and data reduction. The results, including precision, recall, F1 score, and compression ratios for each method, are presented and discussed in the following sections.

\section{Compression based on I/Q Sample Rate and Sample Resolution Reduction}

To reduce the data volume without interpreting the signal content, we first consider two lossy compression strategies: reducing sample resolution and reducing sample rate.

\subsection{Methodology}

\subsubsection{Reducing the Sample Resolution}

Reducing the number of bits per sample is effective only when the resulting
size aligns with full bytes. Given our 12-bit ADC resolution (stored at 16 bits), we consider
reductions to 8 bits and 4 bits per sample. Reducing the resolution to 8 bits halves the data volume, yielding $R = 2$. Further reducing to 4 bits results in another halving as both I and Q can then be packed into a single byte.

The bit-selection strategy for the reduction depends on the expected signal strength:
\begin{itemize}
    \item \textbf{Strong signals:} Drop the least significant bits (LSBs), which
          are likely to represent noise. This is equivalent to performing an
          integer division by a power of two.
    \item \textbf{Weak signals:} Drop the most significant bits (MSBs),
          assuming they are unused (i.e.\ zeroes). Care must be taken to avoid
          signal loss or clipping.
\end{itemize}

We dynamically select the bits to be dropped \emph{per detected signal} based
on actual bits used in I and Q. To determine which bits can be safely discarded, we first estimate the number of bits actually needed to represent the signal's dynamic range. This is done by evaluating the maximum value present in the I and Q components of each detected signal. If the current bit usage exceeds the target resolution, we discard the excess least significant bits (LSBs) by shifting the values to the right. This effectively preserves the most significant structure of the signal while compressing its representation. If the signal already fits within the target bit range, no transformation is applied.

\subsubsection{Reducing the Sample Rate}

Sample rate reduction is achieved via downsampling. This involves applying a low-pass filter to avoid aliasing, followed by retaining every $n$-th sample. In our evaluation, we use an order 8 Chebyshev type I filter for the downconversion. Downsampling by a factor of 2 reduces the sample rate from 12~Msamples/s to 6~Msamples/s and the data volume by 50\%. Hence, the relationship between the downsampling factor and the compression ratio is direct: a downsampling factor of 2 yields a compression ratio $R = 2$. However, this also cuts the signal bandwidth in half, potentially discarding frequency components valuable for signal matching.

% \begin{figure*}[t]
%     \centering
%     \includegraphics[width=0.9\linewidth]{figures/transformation.pdf}
%     \caption{Effects of sample resolution reduction and downsampling on the Mode S interrogation signal.}
%     \label{fig:iq_compression}
% \end{figure*}

%An example of the effects of downsampling and sample resolution reduction on a Mode S interrogation signal from our evaluation data set is shown in \autoref{fig:iq_compression}. \dieter{Since we are exceeding the pages, I would actually remove this figure!}

\subsection{Impact on Accuracy of Correlation-based Matching}

In this section, we evaluate the impact of a reduction in signal rate (downsampling) and lowering the signal resolution on the matching performance. We note that, although we consider them separately in our analysis, both downsampling and sample resolution reduction can be applied together. In that case, the overall compression ratio $R$ is the product of the individual compression ratios. For example, reducing the resolution to 8 bits ($R_1 = 2$) and downsampling by a factor of 4 ($R_2 = 4$) results in an overall compression ratio of $R = R_1 \times R_2 = 8$.

\subsubsection{Sample Resolution}
\label{sec:results_resolution_reduction}

\begin{figure}[t]
    \centering
    \includegraphics[width=\linewidth]{figures/png/resolutions_f1_score.png}
    \caption{F1 score by sample resolution and threshold}
    \label{fig:f1_score_resolutions}
\end{figure}

We evaluated the effects of sample resolution reduction on matching accuracy. Starting from the original 12-bit ADC resolution, we reduced the sample resolution to 8, 4, and 2 bits and computed the F1 scores across different similarity thresholds. The results are shown in \autoref{fig:f1_score_resolutions}.

At full 12-bit resolution, the maximum F1 score of 0.71 is achieved at a similarity threshold of 0.87. This threshold provides the best balance between precision and recall, yielding a recall of 66\% and a precision of 75.8\%. Thus, approximately 34\% of true matches are missed, while 24.2\% of the reported matches are false positives.

We note that this result establishes a baseline for matching accuracy under our specific scenario, including hardware configuration, signal modulation and formats, and channel characteristics. It represents the achievable performance when matching is based solely on signal similarity, without leveraging any prior knowledge of the signal structure, such as preambles, error correction codes, or modulation details.

\begin{figure}[t]
    \centering
    \includegraphics[width=\linewidth]{figures/png/resolutions_lost_bits.png}
    \caption{Effective number of least significant bits lost during resolution reduction}
    \label{fig:resolutions_lost_bits}
\end{figure}

Reducing the resolution from 12 bits to 8 bits has almost no effect on matching performance. The maximum F1 score remains at 0.71, achieved at the same threshold of 0.87. As shown in \autoref{fig:resolutions_lost_bits}, more than 40\% of the signals required no bit shift, and an additional 40\% required removal of only a single least significant bit (LSB). This indicates that most received signals were weak relative to the receiver’s dynamic range, and thus the most significant bits were unused. Consequently, removing these bits causes negligible information loss. Furthermore, the least significant bits of strong signals predominantly contain noise, so their removal has minimal impact on matching performance.

Reducing the resolution to 4 bits introduces noticeable degradation. The maximum F1 score decreases to 0.67, and the optimal similarity threshold shifts slightly lower, to between 0.84 and 0.87. This degradation is attributed to the loss of significant information as more LSBs are discarded.

Further reducing the resolution to 2 bits results in a near-complete loss of usable information. The best achievable F1 score drops sharply to 0.17 at a threshold of 0.82, rendering reliable signal matching infeasible at this extreme compression level.

Overall, the results show that reducing the sample resolution to 8 bits is highly effective, achieving a compression ratio of 2 with negligible loss in matching performance, while more aggressive reductions significantly degrade the system’s ability to match signals accurately.

\subsubsection{Downsampling}
\label{sec:results_downsampling}
\begin{figure}[t]
    \centering
    \includegraphics[width=\linewidth]{figures/png/downsampling_f1_score.png}
    \caption{F1 score by downsampling factor and threshold}
    \label{fig:f1_score_downsampling}
\end{figure}

We evaluated the effects of downsampling on matching accuracy by applying downsampling factors of 2, 4, 8, 16, 32, and 64 to the received signals and computing the resulting F1 scores. The results are shown in \autoref{fig:f1_score_downsampling}.

As described in the previous section, the original I/Q data sampled at 12~MHz (downsampling factor of 1) achieves a maximum F1 score of 0.71 at a similarity threshold of 0.87.

Interestingly, moderate downsampling improves the matching performance. Specifically, downsampling by a factor of 2 increases the maximum F1 score to 0.74 at a threshold of 0.90, and further downsampling by a factor of 4 increases the F1 score to 0.76 at a threshold of 0.92. While this appears counterintuitive at first, since reducing the sample rate usually entails information loss, the improvement can be explained by the characteristics of the signals under consideration. SSR~Mode~S interrogations occupy a significantly narrower bandwidth than the $\pm 6$~MHz captured at the original 12~MHz sampling rate. The effective bandwidth of the signals is around 4~MHz, meaning the GRX3X receiver oversamples the signals substantially, capturing noise from the adjacent bands. Downsampling effectively acts as a low-pass filter, suppressing out-of-band noise and thereby increasing the similarity of true matches, as reflected by the ability to apply stricter (higher) matching thresholds.

The best performing configuration is achieved at a downsampling factor of 4, corresponding to a sampling rate of 3~MHz, with a similarity threshold of 0.92. Under this configuration, the matching recall increases from 66\% to 78\%, substantially reducing the fraction of missed true matches from 34\% to 22\%. The precision slightly decreases from 75.8\% to 73\%, but the overall improvement in recall leads to a higher F1 score.

Further increasing the downsampling factor beyond 4 leads to a significant degradation in matching performance. At high downsampling factors, essential signal characteristics are lost due to the corresponding reduction in effective bandwidth, making reliable matching impossible.

These results demonstrate that, although oversampling can improve timestamp accuracy, for the purpose of similarity-based matching, downsampling to a rate close to the true bandwidth of the underlying signals significantly improves performance by enhancing the signal-to-noise ratio.

\subsection{Summary of the Results}

Overall, under the specific conditions of our study, i.e., the hardware configuration, signal modulation and formats, and channel characteristics observed in our dataset, a combination of reducing the sample resolution to 8 bits and downsampling the sample rate from 12~Msamples/s by a factor of 4 to 3~Msamples/s provides the best trade-off between data volume and matching performance. This configuration achieves a compression ratio of 8 while simultaneously improving the F1 score compared to the uncompressed baseline. Further reductions in resolution or sample rate lead to substantial losses in matching accuracy and are therefore not advisable in this application context.

It is important to note that the similarity threshold used for determining matches significantly impacts the trade-off between precision and recall. Lowering the threshold increases recall, thereby reducing the number of false negatives, but at the cost of a higher false positive rate, which in turn lowers precision. If the central data processor (CDP) can tolerate the increased processing effort resulting from more false matches, this trade-off may be acceptable or even desirable, particularly in applications where minimizing missed detections is critical.

\begin{figure}[t]
    \centering
    \includegraphics[width=\linewidth]{figures/png/recall_vs_precision.png}
    \caption{Precision, recall, and F1 score versus similarity threshold for the best performing resolution reduction (to 8 bits) and downsampling (factor 4) configurations}
    \label{fig:recall_vs_precision}
\end{figure}

\autoref{fig:recall_vs_precision} visualizes how precision, recall, and the F1 score vary with the similarity threshold for the best-performing configuration. The curve illustrates the sensitivity of the system to threshold tuning and provides a means to adjust the matching behavior based on application-specific priorities.

\section{DL-Based Compression with Siamese Networks}
The DL-based compression method follows a two-step optimization process. First, a DL model is trained to learn a signal embedding using the siamese network architecture and contrastive loss function, shown in \autoref{fig:siamese_architecture}. In the second step, this embedding is further compressed by reducing the number of bits required to encode each element in the embedding dimensionality, improving storage and transmission efficiency while preserving essential information.

%%%
\subsection{Methodology}
\subsubsection{Model Architecure}
The deep learning model considered for the compression of the signal is a ResNet architecture, originally designed by O’Shea et al.~\cite{oshea} for modulation classification using raw (IQ) samples. This architecture has demonstrated its effectiveness in learning robust features directly from time-series data, making it well-suited for our task of signal embedding within a Siamese framework.

The ResNet model is characterized by its use of multiple residual blocks, often organized into several residual stacks. These residual connections facilitate the training of deeper networks by mitigating the vanishing gradient problem, allowing the network to learn complex representations. Following the residual stacks, the network typically includes fully connected layers that aggregate the learned features into a final output. In our case, the output of these fully connected layers serves as the lower-dimensional embedding of the input signal.

In our experiments, we consider different input sizes for the ResNet architecture, which directly impacts the number of preceding residual stacks within the network and consequently its computational demands. We also explore a range of embedding dimensionalities, from 4 to 20 with a step of 2, for each considered input size. This allows us to analyze the trade-off between the degree of signal compression and the resulting signal matching accuracy. An overview of the specific model configurations, including input sizes and the corresponding computational complexity in terms of FLOPs, is provided in Table~\ref{tab:model_overview}. Notably, with the exception of the input size and the final embedding dimensionality, all other design parameters of our ResNet architecture, such as the number of filters and kernel sizes within the convolutional layers and residual blocks, are kept consistent with the original design proposed by O’Shea et al.~\cite{oshea}.

The detected signals are normalized to a fixed window size of 256 samples through zero-padding or truncation. For signals exceeding 256 samples, we assume that the initial 256 samples provide sufficient matching information to distinguish between different transmissions. For experiments involving smaller input sizes to the ResNet, we apply a downsampling operation to this 256-sample sequence before feeding it into the network. This allows us to investigate the impact of reduced input data on the performance of our signal matching approach.
\begin{table}[]
    \centering
    \caption{Considered DL model parameters}
    \begin{tabular}{c|c|c|c|c|c}
         Downsample & 1 &2&4&8&16\\\hline
         input size & 256&128&64&32&16\\
         Residual stacks & 6&5&4&3&2\\
         nr.param &420k&406k&392k&379k&365k\\
         Mflops  &6.85&3.54&1.88&1.05&0.64\\
    \end{tabular}
    
    \label{tab:model_overview}
\end{table}



\subsubsection{Model training and threshold determination}
The considered ResNet models are trained based on the Siamese architecture and contrastive learning \cite{chopra2005}. In this learning paradigm, we consider pairs of inputs that either match (come from the same transmission, $Y=0$) or do not match (belong to different transmissions, $Y=1$).

The loss for each pair of signals is computed using the contrastive loss function. This function is designed to learn embeddings where signals from the same transmission (positive pairs) have a small Euclidean distance between their representations ($\mathbf{z}_1$ and $\mathbf{z}_2$), while signals from different transmissions have a distance greater than a margin $m$.

The contrastive loss $\mathcal{L}$ for a single pair is defined as:
\[
\mathcal{L} = (1-Y)\frac{1}{2}(\lVert \mathbf{z}_1 - \mathbf{z}_2 \rVert_2)^2 + Y\frac{1}{2}(\max(0, m - \lVert \mathbf{z}_1 - \mathbf{z}_2 \rVert_2))^2,
\]
Here, $Y$ is an indicator variable: $Y=0$ for signals from the same transmission, and $Y=1$ for signals from different transmissions. $\mathbf{z}_1$ and $\mathbf{z}_2$ are the embeddings produced by the Siamese network for the input pair, and $m$ represents the margin, which we empirically set to 0.7.
All models are trained using 40 epochs with early stopping based on the validation dataset, the adam optimizer and a learning rate of 0.0005.  

\subsubsection{Embedding quantization}
To further compress the learned 32-bit floating-point embeddings from the ResNet model, we employ a uniform quantization scheme. This data-driven approach learns a minimum value ($f_{min}$) and a scaling factor for each feature from the training data to optimally cover the range of the embeddings and reduce the number of bits ($n_{bits}$) required for their representation, potentially with some performance degradation.

The quantization parameters for each feature are derived from the training dataset. For a desired number of bits ($n_{bits}$), the minimum ($f_{min}$) and maximum ($f_{max}$) values are determined along each dimension of the embedding space. The scaling factor for each dimension is then calculated as:
\[
\text{scale} = \frac{f_{max} - f_{min}}{2^{n_{bits}} - 1}
\]
This \texttt{scale} defines the step size between the discrete quantization levels. 
To quantize an embedding, each element of the embedding vector is transformed into a discrete integer level. This transformation involves first subtracting the learned minimum value ($f_{min}$) for that feature, then dividing by the learned scaling factor (\texttt{scale}), and finally rounding the result to the nearest integer. The dequantization process reverses this operation: the integer level is multiplied by the \texttt{scale}, and then the $f_{min}$ is added back to obtain an approximation of the original continuous value. By deriving $f_{min}$ and \texttt{scale} from the training data, the quantization scheme adapts to the distribution of the learned embeddings, aiming to minimize information loss. The number of bits, $n_{bits}$, serves as a hyperparameter that regulates the balance between the compression rate and the representational fidelity of the embeddings.

\subsection{Impact on Accuracy of Deep Learning-Based Matching}
This section analyzes the influence of the downsampling factor, embedding dimensionality, and embedding quantization on the signal matching performance achieved by the proposed deep learning approach.

\subsubsection{Downsampling Factor and Embedding Size}
To evaluate the impact of the downsampling factor and the embedding size on matching accuracy, we first established an optimal decision threshold for each model configuration using the validation dataset. This threshold was determined by maximizing the F1 score, which balances precision and recall. This threshold optimization is considered an integral part of the model training process. The resulting F1 scores for various downsampling factors and embedding dimensions are illustrated in \autoref{fig:f1_score_downsampling_DL}.

\begin{figure}[t]
    \centering
    \includegraphics[width=\linewidth]{figures/png/Downsampling_embedding_green.png}
    \caption{F1 score by downsampling factor and number of features}
    \label{fig:f1_score_downsampling_DL}
\end{figure}

As depicted in Figure \ref{fig:f1_score_downsampling_DL}, the downsampling factor, which is inversely related to the computational cost, significantly affects the maximum achievable F1 score. Notably, a reduction in the downsampling factor from 1 (no downsampling) to 2 results in an approximate decrease of 0.05 in the maximum F1 score. Further investigation reveals that this performance difference is primarily attributed to decreased precision. Specifically, while the recall remains relatively stable around 0.85 for both downsampling factors 1 and 2, the precision drops from approximately 0.67  to 0.60 for a downsampling factor of 2. Employing downsampling factors greater than 2 leads to a substantial degradation in the F1 score, characterized by low values for both recall and precision.

In contrast to the downsampling factor, increasing the embedding dimensionality exhibits a less pronounced impact on the F1 score. The F1 score remains relatively consistent across the tested embedding sizes, with the exception of some minor variations likely attributable to the training process. The highest observed F1 score of 0.77 is achieved for embedding sizes of 8, 14, and 18 when no downsampling is applied. The corresponding recall values for these configurations are 0.87, 0.86, and 0.88, respectively. Given the application in localization, where maximizing the detection of potential signal origins (high recall) is crucial, even if it entails a slight reduction in precision, the model without downsampling and an embedding size of 8 is deemed the most suitable configuration. This choice provides a valuable trade-off between effectively identifying potential matches and maintaining a relatively compact embedding size for efficient processing and transmission.
The results demonstrate that the downsampling factor significantly impacts the matching performance, with higher downsampling leading to a substantial decrease in F1 score, primarily due to reduced precision. In contrast, the embedding size shows a less significant influence on the overall performance.


\subsubsection{Embedding Quantization}
We investigate the impact of embedding quantization on model performance for embedding sizes of 4, 8, 14, and 18 without downsampling. Recognizing the practical constraints of CPU-based processing, our approach involves quantizing each embedding dimension and subsequently packing the resulting bits into bytes. The compressed data volume in bytes to represent a single compressed signal embedding, denoted as $\tilde{V}(S)$, is determined by the following equation:
\begin{equation}
\tilde{V}(S) = \left\lceil\frac{n_{bits} \ k}{8}\right\rceil
\label{eqn:compressed_volum}
\end{equation}
where $n_{bits}$ represents the number of bits used to quantize each of the $k$ embedding dimensions. The total compressed data volume for $N$ signals is then $N \cdot \tilde{V}(S)$. In order to calculate the compression ratio $R$ (as defined in \autoref{eqn:compression_rate}), we need to obtain the total uncompressed data volume. We approximate the total uncompressed data volume by $N \cdot V_{avg}(S)$, where $V_{avg}(S)$ is the average uncompressed data volume per signal, which is equal to  976 bytes for our evaluation dataset.

To establish a baseline for comparison, we consider the unquantized embeddings where each of the $k$ dimensions is represented using a 32-bit floating-point number, requiring 4 bytes per dimension. With embedding sizes of 4, 8, 14, and 18, the initial data volume per embedding is 16, 32, 56, and 72 bytes, respectively. The corresponding initial compression ratios (relative to the  976-byte uncompressed size) are approximately 61, 30.5, 17.5, and 13.5. These represent the compression achieved before any explicit quantization to a lower number of bits is applied in our evaluation. \autoref{fig:data_compress} illustrates the F1 score as a function of the compression ratio for models with embedding sizes 4, 8, 14, and 18.

The highest initial F1 scores ($\sim$0.774) are observed for models with embedding sizes 14 and 18. These models maintain this performance until the compression ratio exceeds approximately 75 and 55, respectively. Beyond these points, further increases in compression lead to a significant drop in F1 score. This performance degradation can be attributed to the reduced number of bits allocated to represent each embedding dimension, occurring around 7 bits per dimension in both cases.

For the model with an embedding size of 8, the initial F1 score is slightly lower ($\sim$0.770). However, it exhibits greater resilience to compression, maintaining a higher F1 score at higher compression ratios compared to the larger embeddings. Performance only begins to decline substantially beyond a compression ratio of approximately 160, again corresponding to roughly 6 bits per embedding dimension.

Finally, the embedding size of 4 presents an interesting trade-off between initial F1 score and compressibility. While starting with the lowest initial F1 score ($\sim$0.75) among the considered models, its performance remains relatively stable up to a very high compression ratio of around 325. This maximum compression also corresponds to approximately 6 bits per embedding dimension, meaning each detection can be compressed to approximately 24 bits or 3 bytes.

Overall, our results demonstrate that all considered models can be compressed down to approximately 6 bits per embedding dimension without a substantial loss in F1 score. This represents a reduction in data volume by a factor of approximately 5.3 compared to the 32-bit floating-point baseline. The highest achievable compression ratio, around 300, is observed for the model with an embedding size of 4.
    
\begin{figure}
    \centering
    \input{figures/tikz/ML\_quantization\_f1}
    \caption{F1 scores for different compression ratios for different embedding sizes}
    \label{fig:data_compress}
\end{figure}


\subsection{Summary of the Results}
The evaluation of our deep learning-based compression and matching approach reveals several key insights into the impact of different design parameters on the overall performance.

Firstly, the downsampling factor significantly influences the signal matching accuracy, primarily affecting the precision of the matches. While minimal downsampling (factor of 1) yields the highest F1 scores, it also incurs the highest computational cost. Increasing the downsampling factor leads to a noticeable degradation in performance, suggesting a trade-off between computational efficiency and matching accuracy. Notably, the embedding dimensionality, within the tested range, exhibits a less substantial impact on the F1 score compared to the downsampling factor. Considering the application's preference for high recall in localization, the model without downsampling and an embedding size of 8 was identified as a favorable configuration, balancing detection probability with a manageable embedding size.

Secondly, embedding quantization presents a substantial opportunity for data compression with minimal impact on matching accuracy, contingent on employing an adequate number of bits per embedding dimension. Considering optimal compression while preserving acceptable performance, an embedding size of 4 quantized to 6 bits per dimension achieves a compression ratio of 325 for a F1 score of 0.75. When prioritizing the highest F1 score, a model with an embedding size of 14 utilizing 7 bits per dimension yields an F1 score of 0.774 at a compression ratio of 75.

Interestingly, smaller embedding sizes, such as 4, demonstrate higher compressibility, achieving compression ratios up to 325 with a relatively stable performance, starting from a slightly lower initial F1 score. 

\section{Comparison of the Proposed Approaches}

\begin{table*}[t]
    \centering
    \caption{Comparison of matching performance and compression rate across different approaches}
    \label{tab:compression_comparison}
    \begin{tabular}{lcccc}
        \toprule
        \textbf{Compression Approach} & \textbf{Compression Ratio} & \textbf{F1 Score} & \textbf{Recall} & \textbf{Precision} \\
        \midrule
        Downsampling (factor 4) & 4.0 & 0.76 & 0.78 & 0.74 \\
        Resolution Reduction (8 bits) & 2.0 & 0.71 & 0.66 & 0.76 \\
        \addlinespace
        Embedding 4, quant. 6 bits & 325.3 & 0.75 & 0.87 & 0.66 \\
        Embedding 8, quant. 6 bits & 162.7 & 0.77 & 0.87 & 0.69 \\
        Embedding 14, quant. 7 bits & 75.1 & 0.77 & 0.86 & 0.71 \\
        \bottomrule
    \end{tabular}
\end{table*}

The resolution reduction and downsampling approaches offer straightforward lossy compression techniques that reduce data volume either in the amplitude domain (by reducing the number of bits per sample) or in the time/frequency domain (by reducing the number of samples). These methods are agnostic to signal type and do not require prior knowledge of signal structure or modulation. As a result, they are broadly applicable across a wide range of frequencies and signal types. Moreover, their computational complexity is low, making them well-suited for implementation on resource-constrained devices. A further advantage is that the compressed signal retains its original time-domain representation, which may support precise time alignment and benefit TDOA estimation at the CDP. The main limitation of these methods is that they do not exploit learned or channel-specific features, which limits the achievable matching performance, especially at higher compression rates.

In contrast, the deep learning-based approach employs a Siamese neural network to extract compact embeddings that preserve only the features most relevant for signal matching. These features are learned directly from data, implicitly capturing certain channel-specific characteristics. As a result, this method can offer improved compression-to-accuracy trade-offs. However, because the extracted features are signal-specific, retraining may be required when the system is deployed in new environments with different signal types or propagation conditions. Furthermore, this approach requires sufficient computational resources at the receiver to generate embeddings in real time.

\autoref{tab:compression_comparison} summarizes the matching performance and compression ratio achieved by the different methods. The results confirm that the deep learning-based approach achieves much higher compression rates at slightly higher  F1 scores, particularly due to strong recall. However, the simplicity and robustness of downsampling and resolution reduction make them attractive for scenarios with strict resource constraints and limited prior knowledge about the signals.

Overall, resolution reduction and downsampling offer simple and efficient solutions suitable for diverse signal types and low-power hardware, whereas the Siamese network-based approach provides higher compression efficiency and recall in environments where training and deployment of neural models are feasible.


% \matthias{TODO: Compare performance results from previous sections, perhaps add table comparing the two approaches}

% \subsection{Discussion}


% \begin{itemize}
%     \item Performance of classical vs. AI-based methods.
%     \item Feasibility of real-time deployment.
% \end{itemize}



% \dieter{Feasibility: Running DL models on the edge are also }
% \dieter{Discussion: The fact that our deep learning-based compression yields a smaller representation than the decoded payload offers a compelling insight: the model appears to learn a representation that is inherently tied to the signal's decoding. This is plausible given the nature of ADS-B Mode S, which utilizes a straightforward PPM modulation where distinct bit values correspond to specific pulse positions, a type of pattern the ResNet architecture is well-suited to recognize. This suggests the model might be learning to represent the underlying information more directly than a bit-by-bit decoding. However, the signal-specific nature of this learning necessitates further investigation to quantitatively determine the generalizability and effectiveness of deep learning compression for signals with unknown or more intricate modulation schemes.}

% \matthias{Discuss different F1 score (like accept more false  matches)}

% \matthias{Future ideas: The signal detection problem was neglected in this work as we focused on the compression. However, detection of signals of unknown modulation and structure poses challenges in the real-world. Simple threshold-based detection (squelch) will generally lead less accurate timestamps compared to when the signals are known and matched filters or similar techniques are used to accurately synchronize with a preamble. Now besides just matching signals at the CDP, our methods could also be extended to support estimating the lag between two detections of the same signal by different receivers that is caused by timestamping inaccuracies. The approaches considered in this work that change resolution or rate of I/Q samples support this naturally as cross-correlation can still be performed on the compressed versions of the signals.
\section{Related Work}
Signal compression techniques have been extensively studied and employed to reduce the data volume while preserving essential information. These methods can be broadly categorized into lossless and lossy compression.

Lossless compression techniques, such as Huffman coding and Lempel-Ziv-Welch (LZW) \cite{sayood_info_theory}, allow for perfect reconstruction of the original data. However, these methods typically achieve limited compression ratios, particularly for continuous-valued signals like IQ data. While critical for applications where no information loss is permissible, their utility may be restricted in scenarios demanding significant data reduction. For instance, Rajendran et al. explored lossless compression in the context of crowdsourced spectrum sensing in their ElectroSense work \cite{rajendran2017electrosense}.

Conversely, lossy compression techniques achieve higher compression ratios by selectively discarding information, aiming to minimize the perceptual or functional impact on the reconstructed signal for the intended application. Common lossy methods for radio signals include downsampling, which reduces the temporal resolution by lowering the sampling rate, and quantization, which decreases the precision of the signal amplitude by using fewer bits per sample \cite{oppenheim_signals}. The selection of downsampling factors and quantization levels directly governs the trade-off between the achieved compression rate and the resulting signal fidelity. These techniques are generally signal-agnostic and relatively straightforward to implement, making them widely adopted as baseline compression strategies.

More advanced lossy compression techniques leverage the statistical properties and inherent redundancies within the signal. Transform coding methods, such as the Discrete Cosine Transform and the Discrete Wavelet Transform, decompose the signal into a new set of coefficients, allowing for the selective removal of less significant components. For example, Zeng et al. implemented the FlexSpec framework \cite{flexspec} based on the Walsh-Hadamard Transform for efficient signal compression, demonstrating the potential of transform-based approaches. These methods can often achieve superior compression ratios compared to simple downsampling and quantization while better preserving perceptually or functionally relevant signal characteristics for specific signal types.

The new paradigm of deep learning and machine learning has shown promising results in the domain of image compression, often achieving higher compression ratios while preserving perceptual quality. While RF signal compression has been exploited for tasks like signal demodulation \cite{demod} and RF signal denoising \cite{denoise}, its application in compression with the goal of signal reconstruction has been relatively limited, with \cite{dlAdaptivecompress} serving as a notable example. More recent research has investigated task-specific learned compression, where the compressed representation is optimized for downstream tasks such as classification or detection \cite{dlcompressRF}. This paradigm shift towards task-aware compression, potentially enabling reduced communication overhead and improved efficiency in distributed systems, motivates our investigation into learning compressed embeddings specifically designed for efficient signal matching in a collaborative localization scenario.

\section{Conclusion and Future Work}

We explored the challenge of enabling accurate signal matching for TDOA-based localization under unknown signal characteristics and bandwidth constraints. Two approaches were investigated: classical lossy compression based on reducing I/Q sample resolution and rate, and a data-driven deep learning method employing a Siamese network to generate compressed embeddings. Our experiments on real-world Mode S radar signals demonstrate that classical compression offers substantial gains at minimal computational cost, while the deep learning approach provides significantly higher compression rates at similar matching performance, achieving compression rates exceeding 300 with F1 scores above 0.75. These results highlight the potential of embedding-based representations in collaborative sensing systems where bandwidth is limited and prior signal knowledge is absent. However, the deep learning-based method requires training on representative data and sufficient compute resources at the sensor, which may limit its deployment in constrained settings.

Future work includes experiments with different channel conditions and signal types, such as those found in the ISM band, as well as the development of optimized embedded implementations of the deep learning approach to improve its applicability in resource-constrained environments.

\begin{acks}
\dieter{As this is a double blind review, we have to use the acks environment :) }
We thank armauisse and the Cyber Defense Campus for the excellent hackathon!

\dieter{I have to check with Sofie/ Hazem, if I should also ack 6G- Bricks!}
This work is supported by the 6G-Bricks project under the EU’s Horizon Europe Research and Innovation Program with Grant Agreement No. 101096954.
\end{acks}
\balance

%%
%% The next two lines define the bibliography style to be used, and
%% the bibliography file.
\bibliographystyle{ACM-Reference-Format}
\bibliography{sample-base}


%%
%% If your work has an appendix, this is the place to put it.
% \appendix

% \section{Research Methods}

% \subsection{Part One}

% Lorem ipsum dolor sit amet, consectetur adipiscing elit. Morbi
% malesuada, quam in pulvinar varius, metus nunc fermentum urna, id
% sollicitudin purus odio sit amet enim. Aliquam ullamcorper eu ipsum
% vel mollis. Curabitur quis dictum nisl. Phasellus vel semper risus, et
% lacinia dolor. Integer ultricies commodo sem nec semper.

% \subsection{Part Two}

% Etiam commodo feugiat nisl pulvinar pellentesque. Etiam auctor sodales
% ligula, non varius nibh pulvinar semper. Suspendisse nec lectus non
% ipsum convallis congue hendrerit vitae sapien. Donec at laoreet
% eros. Vivamus non purus placerat, scelerisque diam eu, cursus
% ante. Etiam aliquam tortor auctor efficitur mattis.

% \section{Online Resources}

% Nam id fermentum dui. Suspendisse sagittis tortor a nulla mollis, in
% pulvinar ex pretium. Sed interdum orci quis metus euismod, et sagittis
% enim maximus. Vestibulum gravida massa ut felis suscipit
% congue. Quisque mattis elit a risus ultrices commodo venenatis eget
% dui. Etiam sagittis eleifend elementum.

% Nam interdum magna at lectus dignissim, ac dignissim lorem
% rhoncus. Maecenas eu arcu ac neque placerat aliquam. Nunc pulvinar
% massa et mattis lacinia.

\end{document}
\endinput
%%
%% End of file `sample-sigconf.tex'.
